% 倍长中线 例3：D 是 BC 中点，AD=4, AB=6, AD⊥AB，求 AC=10
\documentclass[tikz,border=4pt]{standalone}
\usepackage{ctex}
\usepackage{amsmath}
\usetikzlibrary{calc, decorations.markings}
\begin{document}
\begin{tikzpicture}[scale=0.7, thick,
  tick/.style={postaction={decorate, decoration={markings,
    mark=at position 0.5 with {\draw (-1.5pt,-3pt) -- (1.5pt,3pt);}}}}]
  % AD⊥AB, AD=4, AB=6. 设 A=(0,0), B=(-6,0)（AB 向左），AD 向上：D=(0,4)
  % B = (-6,0), D = (0,4), 而 D 是 BC 中点 -> C = 2D-B = (6, 8)
  \coordinate (A) at (0,0);
  \coordinate (B) at (-6,0);
  \coordinate (D) at (0,4);
  \coordinate (C) at (6,8);
  % E: 延长 AD, DE=AD -> E = 2D-A = (0,8)
  \coordinate (E) at (0,8);
  \draw (A)--(B);
  \draw (A)--(C);
  \draw (B)--(C);
  \draw[red, dashed] (A)--(D);
  \draw[red, dashed] (D)--(E);
  \draw[blue, dashed] (C)--(E);
  \draw[tick] (B)--(D);
  \draw[tick] (D)--(C);
  % 直角 AD ⊥ AB
  \draw ($(A)+(0,0.4)$) -- ($(A)+(-0.4,0.4)$) -- ($(A)+(-0.4,0)$);
  \node[below right] at (A) {$A$};
  \node[below] at (B) {$B$};
  \node[above right] at (C) {$C$};
  \node[right] at (D) {$D$};
  \node[above] at (E) {$E$};
  \node[below] at ($(A)!0.5!(B)$) {\footnotesize $6$};
  \node[right] at ($(A)!0.5!(D)$) {\footnotesize $4$};
\end{tikzpicture}
\end{document}
